\section{System Architecture}\label{sec:arch}

\llm is assembled in five stages: (i) from-scratch pre-training of a hybrid language backbone, (ii) agentic post-training (annealing, SFT, RLVR), (iii) from-scratch pre-training of a vision encoder, (iv) multimodal alignment, and (v) semiconductor-domain adaptation.
% 中文：\llm 由五阶段组装：(i) 从零预训练混合语言骨干、(ii) 智能体后训练（退火、SFT、RLVR）、(iii) 从零预训练视觉编码器、(iv) 多模态对齐、(v) 半导体领域适配。
This section fixes what is shared---the two backbones, the serving-state argument, the subagent protocol, the visual-token interface, and the harness contract---before each stage is specified in its own section.
% 中文：本节先固定共享部分——两个骨干、服务状态论证、subagent 协议、视觉 token 接口与框架契约——随后各阶段分节详述。

\subsection{Language Backbone}\label{sec:arch}
The backbone is a sequential hybrid~\cite{dao2024transformers,nvidia2025nemotronh}: 56 layers interleaving Mamba-2 state-space layers, four grouped-query-attention (GQA) layers and MLP layers, reusing Megatron-Core's 56-layer Mamba-hybrid pattern verbatim while halving only the hidden and FFN widths.
% 中文：骨干是顺序式混合模型：56 层交替 Mamba-2、4 层 GQA 与 MLP，原样复用 Megatron-Core 的 56 层 Mamba 混合模式，仅把隐藏维与 FFN 宽度减半。
It is not a Nemotron-H model: an order of magnitude smaller, no mixture-of-experts, pre-trained from scratch on an independent corpus under a fixed GPU-hour budget.
% 中文：它不是 Nemotron-H：小一个数量级、无 MoE、在固定 GPU·h 预算下于独立语料从零预训练。
Each Mamba-2 layer uses 8 state-space groups; the attention layers use 16 query heads and 4 query groups with untied embeddings (Table~\ref{tab:config}).
% 中文：每个 Mamba-2 层 8 个状态空间组；注意力层 16 个查询头、4 个查询组，词嵌入不共享（表 \ref{tab:config}）。
Only the four attention layers carry rotary position embeddings, so long-context extrapolation depends on $4/56=7.1\%$ of layers (Section~\ref{sec:longctx}).
% 中文：只有 4 个注意力层带 RoPE，故长上下文外推仅取决于 4/56=7.1% 的层（第 \ref{sec:longctx} 节）。

\begin{table}[t]
\centering
\caption{Language-backbone configuration (Mamba2-hybrid, $\approx$2.22B).}
\label{tab:config}
\begin{tabular}{lc}
\toprule
\textbf{Hyper-parameter} & \textbf{Value} \\
\midrule
Transformer layers & 56 \\
Architecture & hybrid (Mamba-2 + GQA + MLP) \\
Attention layers & 4 (interleaved) \\
Hidden size & 2048 \\
FFN hidden size & 8192 \\
Attention heads (Q) / groups & 16 / 4 \\
Mamba-2 state-space groups & 8 \\
Vocabulary size & 129{,}281 \\
Sequence length & 16{,}384 (production) \\
Normalisation / activation & RMSNorm / SiLU \\
Parameters & $\approx$2.22B \\
\bottomrule
\end{tabular}
\end{table}

\subsection{Why a Hybrid: Constant-Size Session State}\label{sec:session-state}
An agentic session is not a prompt: each \texttt{run\_code} round appends code, stdout, trace and error text, and the conversation that results is re-read on every subsequent step.
% 中文：智能体会话不是一次提示：每轮 \texttt{run\_code} 都会追加代码、stdout、trace 与报错文本，且后续每一步都要重读整段对话。
For a dense Transformer the cached key/value state grows linearly in the session length, so per-step memory is $O(L)$ and total cost is quadratic in the number of rounds.
% 中文：对稠密 Transformer，缓存的 K/V 状态随会话长度线性增长，故每步内存为 $O(L)$，总成本随轮数呈二次增长。
In \llm only the four GQA layers contribute a growing cache; the 52 Mamba-2 layers carry a recurrent state of fixed size.
% 中文：在 \llm 中只有 4 个 GQA 层贡献增长中的缓存；其余 52 层 Mamba-2 携带固定大小的循环状态。
Writing $L_a{=}4$ attention layers with $h$ KV heads, head dimension $d$, context $L$ and BF16 storage, the session state is
% 中文：记注意力层数 $L_a{=}4$、KV 头数 $h$、头维 $d$、上下文 $L$、BF16 存储，则会话状态为
\begin{equation}
M(L)\;=\;\underbrace{2\,L_a\,h\,d\,L}_{\text{attention cache}}\;+\;\underbrace{C_{\mathrm{ssm}}}_{\text{constant}},
\label{eq:mem}
\end{equation}
where $C_{\mathrm{ssm}}$ is independent of $L$; the corresponding dense model replaces $L_a$ by $56$ and therefore grows $14\times$ faster.
% 中文：其中 $C_{\mathrm{ssm}}$ 与 $L$ 无关；对应的稠密模型把 $L_a$ 换成 56，增长速度快 14 倍。
This is a statement about scaling, not about absolute size, and it is measured, not assumed: Figure~\ref{fig:session-mem} is the placeholder for the measurement we will report.
% 中文：这是关于增长阶而非绝对规模的陈述，且将被实测而非假定：图 \ref{fig:session-mem} 是待填的测量占位。

\begin{figure}[t]
\centering
\phfig{3.1cm}{\textbf{Figure placeholder: session state}\\[2pt]$M(L)$ vs.\ rounds $L$\\hybrid (4 attn) vs.\ dense (56 attn) vs.\ API baseline\\[2pt]\tbd}
\caption{Per-session serving memory against accumulated context. Predicted shape from Eq.~\ref{eq:mem}: the hybrid grows $14\times$ slower in the attention term and adds a constant SSM term. Measurement is pending.}
\label{fig:session-mem}
\end{figure}

\subsection{Subagent--Master Protocol}\label{sec:protocol}
\llm is not meant to replace the orchestrator; it is meant to be one expert per tool, reachable from a stronger planner or from a human flow (Figure~\ref{fig:deploy}).
% 中文：\llm 不是替代编排者，而是「每个工具一位专家」，可由更强的规划器或人工流程调用（图 \ref{fig:deploy}）。
The boundary is deliberately asymmetric: the subagent exchanges \emph{natural language} with its caller and \emph{EDA scripts} with the tool, so a master never has to understand SKILL or Tcl, and the subagent never has to trust the master's code.
% 中文：边界刻意不对称：subagent 与调用者交换自然语言、与工具交换 EDA 脚本——master 无需懂 SKILL/Tcl，subagent 也无需信任 master 的代码。
Table~\ref{tab:protocol} fixes the message schema; every field is machine-checkable, which is what lets us swap the backbone without touching the harness.
% 中文：表 \ref{tab:protocol} 固定消息 schema；每个字段都可机器校验，这正是「换骨干不动框架」的前提。

\begin{table}[t]
\centering
\caption{Subagent--master message schema (T3). $\to$: caller to subagent; $\leftarrow$: subagent to caller.}
\label{tab:protocol}
\scriptsize
\setlength{\tabcolsep}{3pt}
\begin{tabular}{@{}llp{2.55cm}l@{}}
\toprule
\textbf{Field} & \textbf{Type} & \textbf{Constraint} & \textbf{Dir.} \\
\midrule
\texttt{role} & enum & \texttt{layout} / \texttt{schematic} / \texttt{waveform} & $\to$ \\
\texttt{intent} & text & NL request, $\le$512 tok & $\to$ \\
\texttt{artifacts} & image[] & rendered view(s), $224$--$448$px & $\to$ \\
\texttt{budget} & int & max \texttt{run\_code} calls & $\to$ \\
\texttt{script} & code & PyAether / SKILL / Tcl & $\leftarrow$ \\
\texttt{tool\_calls} & record[] & tool id + args (audit) & $\leftarrow$ \\
\texttt{assertion} & record & per-assertion pass/fail & $\leftarrow$ \\
\texttt{handoff} & text & NL summary + open issues & $\leftarrow$ \\
\bottomrule
\end{tabular}
\tnote{Every field is validated at the tool boundary, so the caller cannot inject execution and the subagent cannot claim success without an assertion record.}
\end{table}

\subsection{Visual-Token Interface}\label{sec:vtoken}
The encoder of stage~(iii) renders an image as $14{\times}14{=}196$ tokens at $224$px/$16$-px patches; two coarser-finer tiers are pre-registered so that Claim~B is measured rather than argued (Table~\ref{tab:vtoken}).
% 中文：第 (iii) 阶段的编码器把图像表示为 224px/16px 块下的 196 个 token；另预注册两档更细粒度，使主张 B 可测而非仅争论（表 \ref{tab:vtoken}）。
The tier is a single scalar in the projector, so accuracy and cost move along one axis and the ablation cannot confound architecture with resolution.
% 中文：档位只是投影层里的一个标量，故精度与成本沿同一根轴变化，消融不会把架构与分辨率混为一谈。

\begin{table}[t]
\centering
\caption{Visual-token tiers exposed by the stage-(iv) interface (T4).}
\label{tab:vtoken}
\scriptsize
\setlength{\tabcolsep}{3pt}
\begin{tabular}{@{}lcccc@{}}
\toprule
\textbf{Tier} & \textbf{Input} & \textbf{Tokens} & \textbf{px/token} & \textbf{Rel.\ cost} \\
\midrule
T196 (default) & $224$ & $196$ & $16{\times}16$ & $1.0\times$ \\
T576 & $336$ & $576$ & $16{\times}16$ & $2.9\times$ \\
T1024 & $448$ & $1024$ & $16{\times}16$ & $5.2\times$ \\
\bottomrule
\end{tabular}
\tnote{Rel.\ cost counts visual tokens only; end-to-end serving cost per task is measured separately (Table~\ref{tab:econ}). Which tier is ``enough'' for layout, schematic and waveform views is the object of RQ5--RQ6.}
\end{table}

\subsection{Harness Contract: Byte-Identical}\label{sec:harness}
Because ZhuLong attributes its gains to the interface rather than to the policy, varying the backbone is only informative if everything else is frozen.
% 中文：由于 ZhuLong 把增益归因于接口而非策略，只有在其余全冻结时替换骨干才具信息量。
Table~\ref{tab:harness} is the checklist we commit to; any deviation invalidates the comparison and is reported as such.
% 中文：表 \ref{tab:harness} 是我们承诺的核对清单；任何偏离都会使比较失效，并如实报告。

\begin{table}[t]
\centering
\caption{Harness contract for the \llm--ZhuLong comparison (T5).}
\label{tab:harness}
\scriptsize
\setlength{\tabcolsep}{3pt}
\begin{tabular}{@{}p{3.05cm}p{1.65cm}p{1.35cm}@{}}
\toprule
\textbf{Component} & \textbf{ZhuLong} & \textbf{\llm} \\
\midrule
Agent scaffold & Cline-CLI, unmodified & identical \\
MCP tools & 4 (search, detail, doc, run) & identical \\
Retrieval stack & API + keyword + docs & identical \\
Sandbox / effector & vendor runtime & identical \\
Anti-cheat hook & enabled & identical \\
Assertion criterion & per-task assertions & identical \\
Timeout & $2500$\,s / task & identical \\
Runs & 5, reshuffled, fresh ctx & identical \\
Contamination & $n$-gram blacklist & identical \\
Task schema / slices & $158/157/134$ & identical \\
\bottomrule
\end{tabular}
\tnote{Only the backbone (and, in stage iv, the visual channel) varies. Deviations found during the run are listed in the appendix rather than silently corrected.}
\end{table}

\begin{figure}[t]
\centering
\phfig{2.6cm}{\textbf{Figure placeholder: deployment}\\[2pt]master (planner / human) $\leftrightarrow$ \llm subagent $\leftrightarrow$ EDA runtime\\per-tool expert instances behind one NL contract}
\caption{Deployment shape: one subagent per EDA tool, addressed in natural language, speaking EDA script to the tool.}
\label{fig:deploy}
\end{figure}

